\startbuffer[sacc]
  {\emph Saccharomyces cerevisiae}
\stopbuffer

\startbuffer[bac]
  {\emph Bacillus}
\stopbuffer

\startbuffer[clos]
  {\emph Clostridium}
\stopbuffer

\startchapter[title={Revisão}]
  A industria de biotecnologia é responsável por produções alternativas de compostos usuais na nossa sociedade, utilizando organismos vivos para produzir alimentos, biocombustiveis \cite[authoryear][ref::nielsen_0], farmacos \cite[authoryear][ref::zhang_0] e compostos químicos \cite[authoryear][ref::meadows_0]. No balanço de sustentabilidade, estas rotas alternativas, utilizando organismos modificados ou não, permitem uma aproximação orgânica e amigavel com o meio ambiente para os processos produtivos envolvidos, promovendo uma transição para uma economia verde, ou seja, reduzindo o saldo de carbono emitido \cite[authoryear][ref::naseri_0]. Estas alternativas não são à principio uma degressão em produtividade nas cadeias produtivas mencionadas em prol da sustentabilidade, mas sim, puramente uma evolução em eficiência na alocação de espaço físico, consumo de insumos e geração de subprodutos \cite[authoryear][ref::metcalfe_0]. A sustentabilidade é mais um ponto positivo que pesa a favor de implementações em grande escala de processos que utilizem técnicas em biotecnologia em suas rotas produtivas.

  \startsection[title={Biorreatores}]
    Biorreatores fornecem ambientes favoráveis para o crescimento de culturas de microrganismos e podem ser otimizados para a produção de determinado produto bioquímico, desde que haja acesso a mecanismos de controle das variáveis físicas envolvidas, o que geralmente é oferecido por reatores projetados especificamente para essa finalidade, contrariamente aos bioreatores de ocorrencia natural, onde a manipulaçao consistente de variaveis nao acontece com facilidade ou é completamente mitigada - a exemplo, uma poça de agua a ceu aberto com atividade microbiana em que se deseja controlar a temperatura. A necessidade de um projeto de biorreator decorre da responsabilidade de garantir a reprodutibilidade dos processos bioquímicos desenvolvidos ao longo de uma cadeia produtiva. A possibilidade de controlar algumas variáveis \emdash ou até mesmo todas aquelas avaliadas como impactantes para a conversão final \emdash é uma característica indispensável para aplicações industriais de rotas produtivas desenvolvidas em laboratório. Sendo que uma vez determinadas as condições de operação na etapa experimental, basta ajustar-se as variáveis extensivas do processo para uma escala industrial.

    Segundo \cite[authoryears][ref::erickson_0], existem diversas formas de classificar biorreatores, seja pelo tipo de alimentação ou pelos produtos finais obtidos. Existem três tipos principais de alimentação: batelada, alimentação alternada e alimentação contínua. A alimentação em batelada ocorre quando o substrato é introduzido no sistema apenas no início do processo, sem manipulações externas posteriores na quantidade de fonte de carbono. A alimentação alternada consiste em uma estratégia de distribuição da fonte de carbono seguindo um cronograma, diferenciando-se da primeira pela possibilidade de introdução de substrato durante o experimento. Por fim, a alimentação contínua caracteriza-se pela introdução constante de alimento ao longo de todo o processo, representando uma tentativa de reproduzir um sistema estacionário a partir da manutenção da concentração de substrato no biorreator.

    %\startplacefigure[title={Ilustração de biorreator genérico}]
    %    \externalfigure[image/bioreactor.jpeg][height=200pt]
    %\stopplacefigure

    Quando tratamos de processos aeróbicos, é necessário fornecer oxigênio para o desenvolvimento da cultura. Nessa perspectiva, boa parte dos reatores destinados ao cultivo de organismos aeróbicos necessita de agitação para uniformizar a concentração de oxigênio em todo o meio, o que é naturalmente dificultado pela baixa solubilidade do oxigênio em água. Isso pode ser realizado por meio de reatores de mistura, colunas de ar dissolvido, membranas ou pelo aumento da área de interface líquido-gás, o que não ocorre necessariamente em processos anaeróbicos ou em sistemas envolvendo culturas mistas \cite[authoryear][ref::erickson_0]. Essa última opção tem como proposta estimular a competição entre diferentes variantes presentes com base na seleção natural. Portanto, os tipos de cultura influenciam diretamente no projeto físico do biorreator.

    Outra condição de operação que influencia drasticamente tanto o projeto físico quanto o modelo matemático de um reator é a fase dos reagentes e produtos. Em fase líquida, a cinética bioquímica depende da concentração dos reagentes e, em alguns casos, também da concentração dos produtos. Para muitas reações bioquímicas, essa abordagem cinética pode ser tratada, com boa aproximação, considerando dependência de primeira ordem em baixas concentrações e de ordem zero em concentrações elevadas \cite[authoryear][ref::erickson_0]. Em fase sólida, a cinética bioquímica torna-se mais dependente da disposição espacial do substrato sólido e do contato superficial deste substrato com o microrganismo. Dessa forma, uma solubilidade substancial dos reagentes é essencial para garantir adequada difusão no meio e, consequentemente, a formação do produto de interesse \cite[authoryear][ref::bhargav_0]. Em alguns sistemas, utilizam-se bioenzimas para hidrolisar o substrato em reagentes mais solúveis, aumentando a conversão por unidade de tempo.

    O controle das variáveis físicas envolvidas é importante para a predição da resposta do sistema em relação a qualquer perturbação de entrada, bem como para a manipulação do sistema, de forma que os resultados de réplicas experimentais sejam os mais semelhantes possíveis dentro das limitações de sistemas não determinísticos. Entre as variáveis que podem ser medidas em um biorreator estão temperatura, pressão, pH, potência de agitação, fluxo de gases e líquidos, concentração de gases, concentração de biomassa e tensão interfacial \cite[authoryear][ref::erickson_0].
  \stopsection

  \startsection[title={Fermentação}]
    A fermentação ocupa posição de destaque na indústria biotecnológica, assumindo protagonismo, ainda hoje, em rotas produtivas de compostos essenciais para a sociedade \cite[authoryear][ref::metcalfe_0]. Grande parte das vacinas, medicamentos e biocombustíveis é obtida preferencialmente por processos fermentativos, e sua viabilidade econômica decorre justamente da rota empregada. Caso contrário, métodos alternativos de produção seriam mais trabalhosos, dispendiosos e provavelmente economicamente inviáveis na escala atual.

    Os microrganismos responsáveis pela fermentação podem ser bactérias, leveduras, arqueias, fungos ou células de mamíferos \cite[authoryear][ref::muller_0]. Portanto, existe uma ampla variedade de caminhos metabólicos possíveis considerando o conjunto desses organismos. As leveduras, por exemplo, são responsáveis pela produção de etanol, pão, bebidas alcoólicas, entre outros produtos. reforçando a relevância econômica e social desses microrganismos e de suas aplicações industriais.

    A produção de etanol por fermentação é um dos processos biotecnológicos mais importantes da indústria moderna, sendo amplamente utilizada para a obtenção de biocombustíveis e produtos químicos renováveis \cite[authoryear][ref::dashko_0]. Nesse processo, leveduras como \getbuffer[sacc] convertem açúcares provenientes de matérias-primas vegetais, como cana-de-açúcar, milho e biomassa lignocelulósica, em etanol e dióxido de carbono sob condições anaeróbicas \cite[authoryear][ref::bhargav_0]. A fermentação alcoólica apresenta grande relevância econômica, especialmente em países como o Brasil, onde o etanol é utilizado em larga escala como combustível automotivo. Além de sua aplicação energética, o etanol também é empregado nas indústrias farmacêutica, alimentícia e química, muitas vezes na forma de gel ou liquida em diversas concentracoes para ser utilizado na esterilizacao.

    \startplacefigure[title={Ilustração de levedura}, reference={fig:yeast_ilustration}]
        \externalfigure[image/yeast_cell.jpeg][height=200pt]
    \stopplacefigure

    A \in{Figura}[fig:yeast_ilustration] mostra uma célula de levedura. Nela, é possível observar o núcleo, onde se encontra o material genético principal dos estudos de modificação genética. A mitocôndria também possui material genético, embora a extensão dos estudos relacionados à sua manipulação seja menor quando comparada à do núcleo. Um ponto importante na compreensão das reações químicas de oxirredução \cite[authoryear][ref::muller_0] é que elas ocorrem no citoplasma ou são coordenadas na mitocôndria quando associadas a processos aeróbicos.
  \stopsection

  \startsection[title={Fermentação Alcoólica}]
    No que se refere à produção de álcool a partir de fontes de carbono, as leveduras desempenham papel fundamental. Em especial, \getbuffer[sacc] representa um grupo extremamente relevante nesse processo. Do ponto de vista evolutivo, essas leveduras desenvolveram essa capacidade inicialmente em resposta ao surgimento de frutas, fontes naturais de açúcares simples, possibilitando a conversão desses açúcares em álcool tanto em condições aeróbicas quanto anaeróbicas \cite[authoryear][ref::dashko_0].

    Em geral, o metabolismo de um microrganismo pode ser modificado por meio da inserção de genes reacionais em estequiometria adequada para maximizar determinado produto ou inibir mecanismos metabólicos específicos. \cite[authoryears][ref::meadows_0] trabalhou na modificação do genoma de variantes de \getbuffer[sacc] e demonstrou a possibilidade de aumentar significativamente a seletividade de caminhos metabólicos de interesse, especificamente na produção de precursores de ácidos isoprenoides utilizados no tratamento da malária. Em situações como essa, estudadas por \cite[author][ref::meadows_0], o metabolismo torna-se diferente daquele naturalmente encontrado em variantes de \getbuffer[sacc]. Entretanto, ao restringir a análise ao metabolismo nativo genérico, caracterizado principalmente pelo consumo de glicose e pela formação de etanol, é possível visualizar o mecanismo representado na \in{Figura}[fig:alcoholic_fermentation].

      \startplacefigure[title={Mecanismo genérico parcial de fermentação alcoólica}, reference={fig:alcoholic_fermentation}]
          \externalfigure[image/alcoholic_fermentation.png][height=200pt]
      \stopplacefigure

    A \in{Figura}[fig:alcoholic_fermentation] mostra a esquematização da rota de produção de etanol na fermentação alcoólica, onde os açúcares, especialmente glicose e sacarose, são metabolizados pelas leveduras em condições predominantemente anaeróbicas. Nesse processo, a glicose é convertida em piruvato por glicólise e, posteriormente, em etanol e dióxido de carbono. Sendo que durante a glicólise há formação de \chemical{ATP} e reduçao de \chemical{NADH} (transportador de eletrons no ciclo produtivo), enquanto que na formação de etanol a partir de \chemical{Acetaldeido} o \chemical{NADH} é oxidado, ou seja, um ciclo de oxidação-redução do \chemical{NADH}. Dióxido de carbono é liberado no processo durante a conversão de \chemical{Piruvato} em \chemical{Acetaldeido}.
  \stopsection

  \startsection[title={Fermentação de hidrogênio}]
    Os combustíveis fósseis possibilitaram avanços significativos para o desenvolvimento da sociedade moderna ao longo dos séculos, contribuindo para a consolidação da industrialização, da expansão dos meios de transporte e da intensificação da globalização. Embora a implementação inicial de métodos de geração energética em larga escala tenha sido fundamental para o desenvolvimento socioeconômico da humanidade, grande parte desses modelos apresentam impactos expressivos sobre a manutenção da estabilidade ambiental e da saúde do planeta a longo prazo. Atualmente, considerando o avanço tecnológico e a viabilidade de implementação de alternativas ao modelo energético predominantemente utilizado \cite[authoryear][ref::nanda_0], torna-se relevante destacar os avanços científicos relacionados a métodos alternativos de geração de energia renovaveis. Dentre elas, as rotas de conversão de resíduos orgânicos por meio de tratamentos termoquímicos e bioquímicos destacam-se como algumas das opções mais promissoras em substituição aos combustíveis fósseis \cite[authoryear][ref::nanda_0].

    O hidrogênio possui elevado potencial para utilização em larga escala como combustível. Nesse contexto, pode ser empregado diretamente como um biocombustível gasoso ou convertido em hidrocarbonetos e eletricidade por meio de células a combustível \cite[authoryear][ref::das_0]. Para fins de combustão, o hidrogênio apresenta valor calorífico inferior de aproximadamente \math{120 \unit{MJ} \unit{kg}^{-1}} \cite[authoryear][ref::sarangi_0], equivalente ao fornecido por cerca de \math{2.75 \unit{kg}} de gasolina \cite[authoryear][ref::nanda_0]. Tais características, associadas ao fato de sua combustão gerar água como único produto, tornam o hidrogênio uma alternativa energeticamente vantajosa, promissora e sustentável.

    As rotas para produção de hidrogênio são diversas, podendo ter como matéria-prima hidrocarbonetos, gás natural, água ou biomassa. Entre os principais processos, destacam-se a reforma de gás natural, decomposição térmica do gás natural, gaseificação de carvão, oxidação parcial de hidrocarbonetos, pirólise, eletrólise, fotólise, termólise e as rotas biológicas \cite[authoryear][ref::das_0]. Essas rotas, predominantemente termoquímicas e eletroquímicas, em sua maioria, requerem elevadas quantidades de energia durante sua operação, e são consideradas ambientalmente pouco sustentáveis. Em contraste, a produção biológica de hidrogênio ocorre sob condições próximas à temperatura e pressão ambientes, apresentando, portanto, menor demanda energética ao longo do processo \cite[authoryear][ref::das_0].

    Dentro das possibilidades de obtenção biológica de hidrogênio, encontram-se três rotas alternativas na natureza, cada uma com características que podem viabilizar ou inviabilizar sua aplicação, a depender das condições de operação e dos substratos utilizados: biofotólise da água, fotodecomposição e fermentação de compostos orgânicos \cite[authoryear][ref::das_0]. Essas são as principais alternativas para a biossíntese de hidrogênio, realizadas em ambientes anaeróbicos facultativos ou obrigatórios. Porém, tanto na biofotólise quanto na fotodecomposição, a formação de \chemical{H_{2}} é parcial ou completamente inibida na presença de oxigênio.

    A biofotólise adapta o processo de fotossíntese de plantas e algas para a geração de hidrogênio, em vez da produção de biomassa baseada em carbono \cite[authoryear][ref::das_0]. O processo ocorre por meio da absorção de luz, que promove a quebra da molécula de água em gás hidrogênio e gás oxigênio \cite[authoryear][ref::das_0]. Embora existam nuances relacionadas ao mecanismo completo da reação, de modo geral a formação de hidrogênio ocorre pela ação das enzimas hidrogenase ou nitrogenase, como ilustrado na \in{Figura}[fig:biophotolysis].

    \startplacefigure[title={Ilustração do mecanismo de biofotólise da água}, reference={fig:biophotolysis}]
      \externalfigure[image/water_photolysis.png][height=200pt]
    \stopplacefigure

    Na \in{Figura}[fig:biophotolysis], observa-se que a enzima (\chemical{Hidrogenase} ou \chemical{Nitrogenase}) desempenha um papel central na redução de prótons provenientes da fotólise da água. Os elétrons liberados nesse processo são transferidos até a enzima por meio da proteína \chemical{Ferredoxina}, na sua forma reduzida \chemical{Fd_{red}}. A oxidação da \chemical{Ferredoxina} pela \chemical{Hidrogenase} constitui a etapa final da cadeia de transferência eletrônica, na qual os \chemical{H^{+}} são reduzidos, resultando na formação de gás hidrogênio (\chemical{H_{2}}). Em seguida, a \chemical{Ferredoxina} é regenerada na forma oxidada (\chemical{Fd_{ox}}), sendo posteriormente novamente reduzida a \chemical{Fd_{red}}, o que permite a continuidade do ciclo redox. Observa-se também a formação de oxigênio molecular na etapa inicial da decomposição da água, conforme ilustrado na \in{Figura}[fig:biophotolysis]. Nesse contexto, o processo apresenta dependência de sistemas enzimáticos para viabilizar a redução de prótons a hidrogênio, caracterizando-o como um sistema indireto de produção de \chemical{H_{2}}. Em contrapartida, pode ocorrer a evolução direta de hidrogênio a partir da fotólise da água, configurando um sistema denominado direto \cite[authoryear][ref::sarangi_0]. Ambos os mecanismos podem ocorrer nesse tipo de rota metabólica; enquanto, o sistema direto apresenta maior simplicidade, o sistema indireto depende da atuação de enzimas específicas.

    A fotodecomposição de compostos orgânicos, também denominada fotofermentação, ocorre na presença de bactérias fotossintéticas \cite[authoryear][ref::sarangi_0], como as do genero \getbuffer[bac]. Esse processo caracteriza-se pela degradação de espécies químicas por ação da luz, com produção de hidrogênio molecular, de maneira análoga à biophotólise. A principal diferença entre os processos reside na natureza da fonte de alimentação e, consequentemente, na rota metabólica global envolvida, embora o princípio de conversão fotobiológica seja comparável. No caso da fotodecomposição, a matéria orgânica atua como substrato energético e doador de elétrons, permitindo maior flexibilidade espectral de absorção luminosa \cite[authoryear][ref::das_0]. Associado ao consumo de matéria orgânica, esse processo apresenta a vantagem de possibilitar a utilização de resíduos orgânicos como fonte de entrada, principalmente carboidratos \cite[authoryear][ref::ntaikou_0].

    \startplacefigure[title={Ilustração do mecanismo de fotodecomposiçao}, reference={fig:photo_fermentation}]
      \externalfigure[image/photo_fermentation.png][height=200pt]
    \stopplacefigure

    A \in{Figura}[fig:photo_fermentation] apresenta semelhanças com a \in{Figura}[fig:biophotolysis], destacando-se, em particular, a presença da \chemical{Ferredoxina}, a qual desempenha função equivalente em ambos os mecanismos. Esses dois sistemas diferem principalmente quanto à fonte de substrato, sendo que, no segundo caso, ácidos orgânicos são preferencialmente utilizados em substituição à água como doadores de elétrons. Outra analogia relevante pode ser estabelecida com o sistema apresentado na \in{Figura}[fig:alcoholic_fermentation], no qual o \chemical{NADH} desempenha papel funcional análogo ao da \chemical{Ferredoxina} nos processos de fotodecomposição e biophotólise da água. Em ambos os casos, esses cofatores atuam como transportadores de elétrons, viabilizando reações de redução na estequiometria final dos processos \cite[authoryear][ref::sarangi_0]. A \in{Equaçao}[eq:photo_fermentation] apresenta uma representação estequiométrica genérica do processo de fotodecomposição.

    \startplaceformula[reference={eq:photo_fermentation}]
      \startformula
        \startchemicalformula
          \chemical{2CH_{3}COOH} 
          \chemical{PLUS} 
          \chemical{H_{2}O} 
          \chemical{GIVES} 
          \chemical{CO_{2}} 
          \chemical{PLUS} 
          \chemical{4H_{2}}
        \stopchemicalformula
      \stopformula
    \stopplaceformula

    Na \in{Equação}[eq:photo_fermentation], observa-se a conversão do \chemical{Ácido acético} em \chemical{H_{2}}. Contudo, essa representação não explicita a dependência do processo em relação à enzima \chemical{Nitrogenase}, bem como à necessidade de \chemical{ATP} para sua atividade catalítica \cite[authoryear][ref::muller_0].

    A lista extendida das variações de substrato aceitas por sistemas de fotodecomposição de compostos orgânicos inclui monóxido de carbono em conjunto com água, em um processo denominado translação microbiana \cite[authoryear][ref::das_0], onde a estequiometria correspondente é apresentada na \in{Equação}[eq:shift].

    \startplaceformula[reference={eq:shift}]
      \startformula
        \startchemicalformula
          \chemical{CO} 
          \chemical{PLUS} 
          \chemical{H_{2}O} 
          \chemical{GIVES} 
          \chemical{CO_{2}} 
          \chemical{PLUS} 
          \chemical{H_{2}}
        \stopchemicalformula
      \stopformula
    \stopplaceformula

    A \in{Equação}[eq:shift] solidifica o conceito nomeado translação, ou {\emph water-gas shift}, que se encaixa no contexto da movimentação do oxigênio da água para o monóxido de carbono, oxidando o carbono, reduzindo o hidrogenio, e formando gás hidrogênio.

    Para ambos os sistemas envolvendo fotodecomposição, não há a formação de oxigênio molecular no meio reacional. Dessa forma, o planejamento do sistema é simplificado, uma vez que não há necessidade de controle rigoroso da presença de \chemical{O_2}, o qual pode atuar como um forte inibidor da atividade de hidrogenases. O efeito inibitório do oxigênio está associado à sensibilidade dessas enzimas ao meio oxidante, podendo ocorrer inativação reversível ou irreversível do sítio ativo contendo metais de transição, particularmente centros \chemical{Fe-Fe} ou \chemical{Ni-Fe}. Nessas condições, a presença de \chemical{O_{2}} leva à oxidação desses centros catalíticos, reduzindo significativamente a eficiência da redução de prótons a \chemical{H_{2}} {\bf <CITAÇÃO>}. Em sistemas como a biofotólise da água, essa limitação exige estratégias adicionais de separação temporal ou espacial entre a etapa de geração de oxigênio e a etapa de produção de hidrogênio, a fim de preservar a atividade enzimática.

    A fermentação de hidrogênio constitui o último dos três principais métodos biológicos de produção de \chemical{H_{2}}, apresentando cinética e condições reacionais distintas, o que resulta em características sistemáticas particulares. Bactérias fermentativas podem apresentar elevadas taxas de evolução de hidrogênio, com capacidade de produção tanto na presença quanto na ausência de luz, não havendo dependência de iluminação, o que permite sua atividade contínua em diferentes ciclos ambientais. Além da produção de \chemical{H_{2}} em quantidades relativamente elevadas quando comparadas a outros sistemas biológicos, esses microrganismos geralmente apresentam taxas de crescimento adequadas para aplicações em processos produtivos, o que pode favorecer o aumento da conversão sob condições ambientais favoráveis \cite[authoryear][ref::das_0]. Também denominada fermentação escura, esse processo pode ocorrer em condições estritamente anaeróbias ou facultativamente anaeróbias, dependendo do microrganismo envolvido, com uma estequiometria genérica apresentada na \in{Equação}[eq:dark_fermentation]. Exemplos incluem bactérias do gênero \getbuffer[clos] \cite[authoryear][ref::muller_0].

    \startplaceformula[reference={eq:dark_fermentation}]
      \startformula
        \startchemicalformula
          \chemical{C_{6}H_{12}O_{6}}
          \chemical{PLUS}
          \chemical{2H_{2}O}
          \chemical{GIVES}
          \chemical{2CH_{3}COOH}
          \chemical{PLUS}
          \chemical{CO_{2}}
          \chemical{PLUS}
          \chemical{4H_{2}}
        \stopchemicalformula
      \stopformula
    \stopplaceformula

    A \in{Equação}[eq:dark_fermentation] exemplifica uma das possíveis reações globais realizadas por microrganismos fermentativos anaeróbios heterotróficos, na qual ocorre o consumo de \chemical{Glicose} com formação de \chemical{Ácido acético} e \chemical{H_{2}}. O \chemical{Ácido acético} é um ácido orgânico que pode ser posteriormente consumido por microrganismos fotodecompositores, de modo que a produção global de hidrogênio é potencializada pela combinação entre fermentação escura e fotofermentação \cite[authoryear][ref::muller_0]. Dessa forma, um arranjo em série das reações representadas nas \in{Equação}[eq:dark_fermentation] e \in{Equação}[eq:photo_fermentation] permite maior aproveitamento dos ácidos orgânicos intermediários, resultando em conversão globais mais eficientes de substrato em \chemical{H_{2}}.
  \stopsection
\stopchapter
